\begin{pastexam}[label={px:ncku-107-ee-ds-10}]{成大電機所 107 學年度 問答題第 6 題}
（問答）

\textbf{6. (9 pts) Set}

There exist two sets which have five members and four members, respectively, which are recorded by the
tree structure. Let the number in the root as ID of a set.

\begin{center}
\begin{tikzpicture}[v/.style={draw,circle,minimum size=5mm,inner sep=0pt},font=\small,>={Stealth[length=2mm]}]
  \node[v] (1) at (0,0) {1};
  \node[v] (4) at (-1.1,-1) {4};
  \node[v] (2) at (-0.4,-1) {2};
  \node[v] (3) at (0.25,-1) {3};
  \node[v] (5) at (1.0,-1) {5};
  \foreach \c in {4,2,3,5} { \draw[->] (\c) -- (1); }
  \node[v] (6) at (2.6,0.2) {6};
  \node[v] (7) at (2.15,-1) {7};
  \node[v] (0) at (3.3,-1) {0};
  \node[v] (8) at (2.7,-2) {8};
  \draw[->] (7) -- (6); \draw[->] (0) -- (6); \draw[->] (8) -- (7);
  \node[font=\small\sffamily] at (-0.1,-2.7) {Set \{1, 4, 2, 3, 5\}};
  \node[font=\small\sffamily] at (2.9,-2.7) {Set \{6, 7, 0, 8\}};
\end{tikzpicture}
\end{center}
\begin{enumerate}[label=\alph*.]
\item \textbf{(6 pts)} Union is an operation to merge two sets. Please respectively show the resulting
data structures after the function WeightUnion1(1, 6) is called when you use number of nodes as the
weight and when you use height as the weight.
\item \textbf{(3 pts)} Find is an operation to find an ID which contains a member. Please show the
resulting data structure if the function \textit{CollaspingFind}(8) is called in the second set, where
collapse aims to reduce the height of a tree.
\end{enumerate}
\end{pastexam}

\begin{pastsoln}
集合用樹表示，每個節點指向父節點，根代表整個集合（資結書第~13.2~節）。 先整理兩棵樹：
\begin{center}
\small
\begin{tabular}{lccc}
\toprule
集合 & 根 & 節點數 & 高度（層數） \\
\midrule
$\{1, 4, 2, 3, 5\}$ & $1$ & $5$ & $2$（$1$ 底下直接是 $4, 2, 3, 5$） \\
$\{6, 7, 0, 8\}$ & $6$ & $4$ & $3$（$8 \to 7 \to 6$） \\
\bottomrule
\end{tabular}
\end{center}
加權合併的規則是「輕的樹的根接到重的樹的根底下」，只改一個指標。 兩種權重剛好給出相反的結果，這就是這題
的重點。

\textbf{a.} \textbf{用節點數當權重}（依大小合併，資結書虛擬碼~13.3）：$5 > 4$，所以根 $6$ 接到根 $1$ 底下，$p[6] = 1$。
新集合的 ID 是 $1$，大小 $9$，高度 $4$（$8 \to 7 \to 6 \to 1$）。

\textbf{用高度當權重}（依高度合併，資結書第~13.2.2~節的秩在不做路徑壓縮時就是高度）：$2 < 3$，所以根 $1$ 接到根 $6$
底下，$p[1] = 6$。 新集合的 ID 是 $6$，高度仍是 $3$（$1$ 那棵掛在深度 $1$，它的子節點在深度 $2$，和 $8$ 一樣深）。
\begin{center}
\begin{tikzpicture}[v/.style={draw,circle,minimum size=5mm,inner sep=0pt},font=\small,>={Stealth[length=2mm]}]
\begin{scope}
  \node[v] (1) at (0,0) {1};
  \node[v] (4) at (-2.0,-1) {4};
  \node[v] (2) at (-1.3,-1) {2};
  \node[v] (3) at (-0.6,-1) {3};
  \node[v] (5) at (0.1,-1) {5};
  \node[v] (6) at (1.3,-1) {6};
  \node[v] (7) at (0.9,-2) {7};
  \node[v] (0) at (1.7,-2) {0};
  \node[v] (8) at (0.9,-3) {8};
  \foreach \c in {4,2,3,5,6} { \draw[->] (\c) -- (1); }
  \draw[->] (7) -- (6); \draw[->] (0) -- (6); \draw[->] (8) -- (7);
  \node at (-0.2,-3.8) {依節點數合併};
\end{scope}
\begin{scope}[xshift=6cm]
  \node[v] (6) at (0,0) {6};
  \node[v] (7) at (-1.2,-1) {7};
  \node[v] (0) at (-0.3,-1) {0};
  \node[v] (1) at (1.2,-1) {1};
  \node[v] (8) at (-1.2,-2) {8};
  \node[v] (4) at (0.15,-2) {4};
  \node[v] (2) at (0.85,-2) {2};
  \node[v] (3) at (1.55,-2) {3};
  \node[v] (5) at (2.25,-2) {5};
  \draw[->] (7) -- (6); \draw[->] (0) -- (6); \draw[->] (1) -- (6); \draw[->] (8) -- (7);
  \foreach \c in {4,2,3,5} { \draw[->] (\c) -- (1); }
  \node at (0.3,-3.8) {依高度合併};
\end{scope}
\end{tikzpicture}
\end{center}
用父節點陣列表示（下標 $0$ 到 $8$，根記為 $-1$）：
\begin{center}
\small
\begin{tabular}{l*{9}{c}}
\toprule
$i$ & 0 & 1 & 2 & 3 & 4 & 5 & 6 & 7 & 8 \\
\midrule
合併前 $p[i]$ & 6 & $-1$ & 1 & 1 & 1 & 1 & $-1$ & 6 & 7 \\
依節點數 & 6 & $-1$ & 1 & 1 & 1 & 1 & 1 & 6 & 7 \\
依高度 & 6 & 6 & 1 & 1 & 1 & 1 & $-1$ & 6 & 7 \\
\bottomrule
\end{tabular}
\end{center}
（Horowitz 的寫法把根的 $p$ 存成 $-(\text{節點數})$，依節點數合併後是 $p[1] = -9$，其他相同。） 依節點數合併時，
樹變高了一層；但這個規則保證的是「$n$ 個節點的樹高度不超過 $\lfloor \lg n \rfloor + 1$」（資結書定理~13.2），這裡
$4 \le \lfloor \lg 9 \rfloor + 1 = 4$，沒有違反。

\textbf{b.} 照題目，在第二個集合（合併前的 $\{6, 7, 0, 8\}$）呼叫 \textit{CollapsingFind}$(8)$（原題拼成 Collasping；就是資結書虛擬碼~13.5 的
路徑壓縮）：
\begin{enumerate}
\item 第一趟從 $8$ 往上找根：$8 \to 7 \to 6$，$6$ 是根，回傳 ID $6$。
\item 第二趟把路上經過的每個節點直接接到根：$p[8] \gets 6$；$7$ 本來就指向 $6$，不變。
\end{enumerate}
結果 $6$ 底下直接是 $7$、$0$、$8$，高度從 $3$ 降成 $2$；之後再找 $8$ 只要一步。
\begin{center}
\begin{tikzpicture}[v/.style={draw,circle,minimum size=5mm,inner sep=0pt},font=\small,>={Stealth[length=2mm]}]
  \node[v] (6) at (0,0) {6};
  \node[v] (7) at (-0.9,-1) {7};
  \node[v] (0) at (0,-1) {0};
  \node[v] (8) at (0.9,-1) {8};
  \draw[->] (7) -- (6); \draw[->] (0) -- (6); \draw[->] (8) -- (6);
\end{tikzpicture}
\end{center}
若是在 a 依節點數合併之後才呼叫，路徑是 $8 \to 7 \to 6 \to 1$，壓縮後 $8$、$7$ 都直接指向 $1$（$6$ 本來就是），回傳 $1$。
\end{pastsoln}
